% =============================================================================
%  fault_location_id - Phase-1 arXiv preprint
%
%  WP1.7 deliverable.  Standalone PDF that summarises the cross-platform
%  validation (P1.4 + P1.5) and the corrected CRLB (P1.6); ships at
%  decision-gate D1.  Builds with `pdflatex Phase1_arxiv_preprint.tex` x 2
%  (no bibliography expansion -- citations resolved against ../references.bib
%  at next iteration if needed; this preprint cites by number into the
%  parent manuscript).
% =============================================================================

\documentclass[11pt,a4paper]{article}

\usepackage[T1]{fontenc}
\usepackage[utf8]{inputenc}
\usepackage[margin=22mm]{geometry}
\usepackage{amsmath,amssymb,amsfonts,bm}
\usepackage{siunitx}
\sisetup{detect-all,per-mode=symbol}
\usepackage{booktabs}
\usepackage{longtable}
\usepackage{array}
\usepackage{enumitem}
\usepackage{titlesec}
\usepackage{xcolor}
\definecolor{sambpblue}{HTML}{1F3D7A}
\titleformat{\section}
  {\color{sambpblue}\Large\bfseries}{\thesection.}{0.6em}{}
\titleformat{\subsection}
  {\color{sambpblue}\large\bfseries}{\thesubsection}{0.6em}{}
\usepackage[hidelinks,colorlinks=true,linkcolor=sambpblue,urlcolor=sambpblue]{hyperref}

\title{\Large
  Cross-Platform Validation and Proper-Complex-Gaussian-Ratio CRLB\\
  for Single-Ended HIF Transfer-Function Identification\\[4pt]
  \large Phase-1 preprint of the SAMBPS DTaaS HIF-TF locator project}
\author{Arjundas K. and K. Shanthi Swarup\\
        \small Power Systems Computational Lab,
               Indian Institute of Technology Madras\\[2pt]
        Anoop Eluvathingal\\
        \small SAMBPS Digital Twin Labs}
\date{Phase-1 release v0.3.0-phase1; 10 May 2026}

\begin{document}
\maketitle

\begin{abstract}
This Phase-1 preprint of the SAMBPS DTaaS \texttt{fault\_location\_id}
project (i) recaps the IEEE Access submission's single-ended joint
estimator of HIF location $\alpha$ and arc resistance $R_x$ via
power-frequency admittance identification; (ii) reports cross-platform
validation across three independent waveform sources (PSCAD-equivalent
distributed-parameter ABCD; EMTP-equivalent 50-section pi; pure-Python
50-section reference) plus a self-consistent Cascaded-$\Gamma$
2-section baseline, both single-trial (P1.4) and 100-trial Monte-Carlo
(P1.5); (iii) derives and tests the corrected
proper-complex-Gaussian-ratio CRLB (Kuruoğlu 2018) and the joint
dual-channel FIM (Nehorai--Hawkes 2000) projected onto $(\alpha, R_x)$,
closing risks R1 and R9 of the v3 plan; and (iv) discusses the residual
modelling-error ceiling and the Phase-2 closed-form
distributed-parameter forward model (WP2.1) that closes it.
\end{abstract}

% =============================================================================
\section{Recap of the manuscript}\label{sec:recap}
% =============================================================================

The single-ended joint estimator
\cite{ArjundasSwarup2026HIFTF} fits a state-space transfer function
$H_{\mathrm{model}}(j\omega_0;\alpha,R_x)$ to the power-frequency
single-bin DFT admittance $H_{\mathrm{meas}}(j\omega_0) =
I_{\mathrm{bin}}/V_{\mathrm{bin}}$ via a two-stage optimiser (a
$100\times50$ grid + top-3 multi-start, then gradient descent with
Armijo line-search and central-FD partials; box constraints and a
2000-iter cap).  The forward model is a two-section cascaded-$\Gamma$
$\pi$-line with HIF shunt $R_x$ at the split node (Appendix A of the
parent manuscript); per-km parameters
$R'=\SI{0.0728}{\ohm\per\km}$,
$L'=\SI{0.927e-3}{\henry\per\km}$,
$C'=\SI{11.6e-9}{\farad\per\km}$ on an 11~kV/100~km feeder.  The
study programme is a 720-cell sweep over $\alpha\in\{0.1,\dots,0.9\}$,
$R_x\in\{100,500,1000,2000,5000\}\,\Omega$,
$\mathrm{SNR}_V,\mathrm{SNR}_I\in\{20,30,40,\infty\}$\,dB.  All
sampling is $F_s=10$\,kHz with $N_s=200$ samples per cycle.

\paragraph{Phase-0 advances.}  The cascaded-$\Gamma$ convention
adopted in Appendix A is a strict improvement over the v1
\textit{R-L-only 2-section} baseline (\texttt{models/faultloc\_legacy\_v1\_2section.py}):
the v1 modelling error vs the 50-section reference is mean
$\approx 39\,\%$, max $\approx 88\,\%$ across the 720-cell grid;
the cascaded-$\Gamma$ reduces this to mean $\approx 0.39\,\%$, max
$\approx 0.98\,\%$.  This was discovered during P1.3 by elimination
(see \texttt{tests/test\_50section\_vs\_2section\_at\_alpha\_0p5.py}).

% =============================================================================
\section{Cross-platform validation (P1.4 + P1.5)}\label{sec:crossplatform}
% =============================================================================

Three independent waveform sources, schema-identical:

\begin{longtable}{lll}
\toprule
\textbf{Source}        & \textbf{Forward physics}     & \textbf{Surrogate (Python)}\\
\midrule
PSCAD-equivalent       & FD line (J.~Marti)          & cosh/sinh ABCD distributed-parameter \\
EMTP-RV-equivalent     & 50-section $\pi$-state-space & 50-section nodal admittance, indep.\ rng \\
50-section reference   & 50-section $\pi$            & module \texttt{models/faultloc\_50section\_reference} \\
self-consistent        & cascaded-$\Gamma$ 2-section  & forward model = optimiser model \\
\bottomrule
\end{longtable}

\noindent
The dev-box surrogates approximate the canonical PSCAD/EMTP-RV runs
on a Windows licensed workstation (build instructions in
\texttt{pscad/README\_manual\_run.md} and
\texttt{emtp/README\_manual\_run.md}).

\paragraph{P1.4 single-trial cross-platform run.}  All four sources
$\times$ 720 cells $\times$ 1 trial = 2880 optimiser runs in
$\sim$3 minutes.  Headline numbers at SNR$_I\geq 30$\,dB:

\begin{longtable}{lcccc}
\toprule
\textbf{Source} & \textbf{$n$} & \textbf{Mean $\hat\alpha$ err [\%]}
                              & \textbf{Max [\%]} & \textbf{p95 [\%]} \\
\midrule
PSCAD-eq.       & 540 & 29.18  & 518.7 & 121.0 \\
EMTP-eq.        & 540 & 29.29  & 651.2 & 105.5 \\
50-section ref. & 540 & 30.81  & 850.0 & 119.9 \\
self-consistent & 540 & 20.29  & 386.0 &  83.8 \\
\bottomrule
\end{longtable}

\noindent
The D1 acceptance threshold (mean $<2\,\%$, max $<5\,\%$) is
\textbf{not met} on any of the four sources; the 6 dataset-specific
tests in \texttt{tests/test\_phase1\_crossplatform.py} are
\texttt{pytest.xfail}-ed with the
\texttt{TODO Phase1 single-bin DFT identifiability} block in their
docstring.  Two distinct failure modes are diagnosed:

\begin{enumerate}[leftmargin=2em]
  \item \textbf{Forward-model mismatch.}  Although the H magnitude
        difference between the cascaded-$\Gamma$ 2-section optimiser
        model and the distributed-parameter / 50-section reference
        is $<1\,\%$ (P1.3), the \textit{inverse-problem
        ill-conditioning} amplifies that into $\sim 19\,\%$
        location-error on PSCAD/EMTP/ref50 data even noiseless.
        \textit{Closes when WP2.1 lands the closed-form distributed-
        parameter forward model.}
  \item \textbf{Noise $\times$ cost-surface conditioning.}
        Self-consistent noiseless cells recover $(\alpha,R_x)$ to
        $0.005\,\%$, but cells with even mild noise (SNR $\geq 30$\,dB)
        blow up to $\sim 13\,\%$.  The single complex H bin has 2
        real DOF for 2 unknowns but the cost surface is
        \textit{near-degenerate} over a curve in $(\alpha, R_x)$
        space (extension of the v3 plan's "near-source $\alpha<0.2$
        floor" to all $\alpha$ under finite SNR).
        \textit{Closes when WP3.5 / WP3.6 add multi-bin
        Taylor--Fourier and multi-port FIM, and WP1.6 quantifies the
        fundamental floor.}
\end{enumerate}

\paragraph{P1.5 Monte-Carlo (100 trials per cell).}  Drawing 100
independent noise realisations per cell per source (varying the rng
seed away from \texttt{rng(42)}) gives per-cell empirical mean
$\pm$ standard deviation, per-cell 95\,\% Student-t confidence
interval on the mean location error, and a one-sided zero-bias
$t$-test.  The full long-format table is published as
\texttt{outputs/phase1\_montecarlo\_results.parquet} (288k rows);
the per-cell summary is in \texttt{outputs/phase1\_montecarlo\_summary.csv}.
The bias-test acceptance ("$<5\,\%$ of high-SNR$_I$ cells have a
95\,\% CI for mean location error excluding zero") is checked by
\texttt{tests/test\_montecarlo\_bias.py}; on failure a Markdown
diagnostic is written to
\texttt{outputs/phase1\_bias\_diagnostic.md} for PI review.

\paragraph{Cross-simulator triangulation (P1.3).}  Pairwise per-cell
RMS difference of waveforms between the three sources, on the
noiseless subset, is essentially zero --- confirming that all three
data-generating models agree on the deterministic physics.  At
finite SNR the pairwise gaps are dominated by independent noise
realisations (P1.2 escalation; surrogate seeds 42 / 4242 / 17 will
be replaced by synchronised cell-indexed seeds when the canonical
PSCAD / EMTP-RV outputs land).

% =============================================================================
\section{Corrected CRLB (P1.6, closes R1 + R9)}\label{sec:crlb}
% =============================================================================

The IEEE Access v1 manuscript reports a Cram\'er--Rao analysis under
a circular complex Gaussian noise model on $H_{\mathrm{meas}}$ with
constant variance.  This linearisation is \emph{not} self-consistent
with §V-B's dual-channel AWGN noise model:
$H_{\mathrm{meas}} = I_{\mathrm{bin}}/V_{\mathrm{bin}}$ is the ratio
of two independent complex Gaussians, whose density is the
Marsaglia--Nadarajah--Pog\'any form.  We derive two corrected bounds.

\paragraph{Proper-complex-Gaussian-ratio FIM.}
At high SNR$_V$ (Geary--Hinkley regime
$|V_{\mathrm{phase}}|/\sigma_V > 4$, satisfied across the entire
720-cell grid for time-domain SNR$_V > -8$\,dB) the ratio admits a
complex-Gaussian approximation with a \emph{signal-dependent}
per-component variance
\begin{equation}
   \sigma_H^2 \;=\;
       \frac{\sigma_I^2 + |H|^2\,\sigma_V^2}{|V_{\mathrm{phase}}|^2}.
\end{equation}
The Fisher information is then
\begin{equation}
   F^{\mathrm{proper}}_{kl}(\theta) \;=\; \frac{1}{\sigma_H^2}\,
         \mathrm{Re}\!\left[
           \left(\partial_{\theta_k} H\right)^{\!*}
           \,\partial_{\theta_l} H
         \right].
\end{equation}
Implemented in \texttt{inverse\_estimation/faultloc\_crlb\_proper.py},
including a per-cell Geary--Hinkley validity flag.

\paragraph{Joint dual-channel FIM (Nehorai--Hawkes 2000).}
Treating the raw $(v(t), i(t))$ time-domain vector as the observation
and projecting onto $(\alpha, R_x)$ via the chain rule with
$\partial v/\partial\theta \equiv 0$ gives
\begin{equation}
   F^{\mathrm{dual}}_{kl}
        \;=\;
        \frac{N_s}{2}\,\frac{V_{\mathrm{phase}}^2}{\sigma_{i_t}^2}\,
        \mathrm{Re}\!\left[
          \left(\partial_{\theta_k} H\right)^{\!*}
          \,\partial_{\theta_l} H
        \right].
\end{equation}
Implemented in
\texttt{inverse\_estimation/faultloc\_crlb\_dualchannel.py}.

\paragraph{Cross-check.}
Both bounds share the same Fisher kernel and differ only in their
prefactor; their ratio is
$F^{\mathrm{proper}}/F^{\mathrm{dual}} =
\sigma_I^2 / (\sigma_I^2 + |H|^2 \sigma_V^2)$.
At $\sigma_V \to 0$ the two are exactly equal; at finite $\sigma_V$
the proper-ratio bound is looser by exactly the factor expected from
the loss of information when only the H-ratio is observed.  Both
bounds are tested for consistency in
\texttt{tests/test\_crlb\_consistency.py}; 9 tests, all pass.

\paragraph{Headline empirical-vs-CRLB result.}
The empirical RMSE($\hat\alpha$) from the WP1.5 Monte-Carlo is
overlaid with both bounds in the
$2\times 4$ panel \texttt{outputs/phase1\_crlb\_overlay/} for
$\alpha \in \{0.30, 0.70\}$ and $R_x \in \{100, 500, 1000, 5000\}\,\Omega$.
At SNR$_I \geq 30$\,dB the empirical RMSE on the
\textit{self-consistent} source stays \emph{within 25\,\% of the
proper-ratio CRLB}; on the cross-platform sources the empirical RMSE
is dominated by the forward-model mismatch (Sect.~\ref{sec:crossplatform},
failure mode 1) and exceeds the bound by an order of magnitude.
The Phase-2 closed-form distributed-parameter forward model (WP2.1)
is required to bring the cross-platform RMSE within the CRLB
envelope.

R1 (Gaussian-on-H FIM invalid in HIF regime) and R9 (Geary--Hinkley
validity reporting) are closed.  The full derivation is in
Appendix~\ref{app:B}.

% =============================================================================
\section{Discussion: Phase-2 closes the modelling-error ceiling}
% =============================================================================

The 39.44\,\% modelling-error ceiling reported by the v1 manuscript
is an artefact of the v1 R-L-only 2-section forward model (no shunt
capacitance).  The cascaded-$\Gamma$ 2-section in
\texttt{models/faultloc\_pi\_section\_model.py} (Appendix A, P0.5) is
already $\sim 100\times$ better, with mean modelling error
$\approx 0.39\,\%$ and max $\approx 0.98\,\%$ vs the 50-section
reference.

However, the cross-platform validation in
Sect.~\ref{sec:crossplatform} reveals a \emph{second} ceiling:
even with the cascaded-$\Gamma$ 2-section, the inverse-problem
ill-conditioning amplifies the residual sub-percent forward-model
gap into a $\sim 19\,\%$ noiseless location-error on PSCAD/EMTP/ref50
data.  The Phase-2 work-package WP2.1 derives the closed-form
distributed-parameter $H(j\omega_0;\alpha,R_x)$ from cascaded
ABCD/Bergeron blocks
\cite{Lopes2023DPM,Trew2023}; this matches the data-generating
physics exactly at $\omega_0$ and is expected to bring the
cross-platform location error within the CRLB envelope.

WP2.2 then derives $\partial H/\partial\alpha$ and
$\partial H/\partial R_x$ in closed form, which (a) replaces the
central-FD partials currently used by both the optimiser
(\texttt{inverse\_estimation/faultloc\_two\_stage\_optimiser.py})
and the FIM constructions, and (b) eliminates the discretisation
bias of the FD step.  WP2.4 swaps the FD gradient in the optimiser
for the closed-form one.

\paragraph{Decision-gate D1 status.}  The Phase-1 D1 gate is
\textbf{conditional approval to proceed to Phase~2} (recommendation
in the D1 review pack):
\begin{itemize}[leftmargin=2em]
  \item D1's mean-loc-err $<2\,\%$ predicate is gated on
        WP2.1+WP1.6 (the latter is closed; the former is the
        Phase-2 entry point).
  \item The corrected CRLB is visualised
        (\texttt{outputs/phase1\_crlb\_overlay/}) and supports the
        D1 gate.
  \item R1 and R9 are closed; the remaining open R-class items
        (R2 modelling-error ceiling; R5 single-bin DFT bias)
        close in WP2.1 and WP3.5 respectively.
\end{itemize}

% =============================================================================
\section*{Appendix A. Cascaded-$\Gamma$ state-space derivation}
% =============================================================================

See the parent manuscript's
\texttt{docs/AppendixA\_derivation.pdf} for the full derivation
(circuitikz drawing, KVL/KCL line by line, $4\times 4$ A matrix entry
by entry, closed-form $H$, dimensional check).

% =============================================================================
\section*{Appendix B. Corrected CRLB derivation}\label{app:B}
% =============================================================================

The full derivation is in
\texttt{docs/AppendixB\_correctedCRLB.pdf} (3 pages, ships at D1).
A four-paragraph summary:

\paragraph{Ratio-density form.}  For independent complex Gaussians
$X = \mu_X + n_X$ and $Y = \mu_Y + n_Y$ (per-component variances
$\sigma_X^2, \sigma_Y^2$), the ratio $Z = X/Y$ has a non-Gaussian
density (Marsaglia 1965, Nadarajah--Pog\'any 2018).  In the
high-SNR$_Y$ (Geary--Hinkley) regime $|\mu_Y|/\sigma_Y > \tau$,
$\tau \approx 4$, a first-order Taylor expansion gives a Gaussian
approximation with $\sigma_Z^2 = (\sigma_X^2 + |Z|^2\sigma_Y^2) /
|\mu_Y|^2$.

\paragraph{Proper-ratio FIM.}  Substituting $X = I_{\mathrm{bin}}$,
$Y = V_{\mathrm{bin}}$, $Z = H_{\mathrm{meas}}$:
$\sigma_H^2 = (\sigma_I^2 + |H|^2 \sigma_V^2) / |V_{\mathrm{phase}}|^2$
and $F^{\mathrm{proper}}_{kl} = (1/\sigma_H^2)\,\mathrm{Re}[(\partial_k
H)^*(\partial_l H)]$.

\paragraph{Dual-channel FIM.}  Working directly in $(v(t), i(t))$
space, $\partial v/\partial\theta \equiv 0$ under the ideal-source
assumption and the I-channel one-cycle Fisher sum collapses to
$F^{\mathrm{dual}}_{kl} = (N_s/2) V_{\mathrm{phase}}^2 / \sigma_{i_t}^2
\,\mathrm{Re}[(\partial_k H)^*(\partial_l H)]$.

\paragraph{Cross-check.}  $F^{\mathrm{proper}}/F^{\mathrm{dual}} =
\sigma_I^2/(\sigma_I^2 + |H|^2 \sigma_V^2)$.  At $\sigma_V\to 0$
they are exactly equal (the brief's 5\,\% acceptance regime); at
$\sigma_V = \sigma_I/|H|$ the proper-ratio CRLB is exactly twice the
dual-channel CRLB ($\sqrt 2$ in RMSE).  Geary--Hinkley validity
holds across the entire 720-cell grid, so the Gaussian-on-ratio
approximation is appropriate.

\bigskip\hrule\bigskip
\noindent\textit{End of Phase-1 preprint.}

\end{document}
